\subsection{权与本原元}

设 V 是一个 g-模。对一切 \(\lambda\in h^{*}\) ，用 \(V^{\lambda}\) 表示把 V 视为 h-模时相对于 \(\lambda\) 的准素子空间（第七章 §1 第 1 节）。\(V^{\lambda}\) 的元素称为 g-模 V 的权为 \(\lambda\) 的元素。诸 \(V^{\lambda}\) 的和是直和（第七章 §1 第 1 节 命题 3）。对一切 \(\alpha\in h^{*}\) 与 \(\lambda\in h^{*}\) ，\(g^{\alpha}V^{\lambda}\subset V^{\alpha+\lambda}\) （第七章 §1 第 3 节 命题 10 (ii)）。\(V^{\lambda}\) 的维数称为 \(\lambda\) 在 V 中的重数；若它 \(\geq1\) ，即若 \(V^{\lambda}\neq0\) ，则称 \(\lambda\) 是 V 的一个权。若 V 是有限维的，由 h 的元素定义的 V 的位似是半单的，故 \(V^{\lambda}\) 是由这样的 \(x\in V\) 组成的集合：\(Hx=\lambda(H)x\) 对一切 \(H\in h\) 成立。

引理 1. 设 \(\mathrm{V}\) 是一个 \(\mathfrak{g}\) -模，\(v \in \mathrm{V}\) 。下列条件等价：

\begin{enumerate}
\def\labelenumi{(\roman{enumi})}
\item
  \(\mathfrak{b}_{+}v\subset kv;\)
\item
  \(\mathfrak{h}v \subset kv\) 且 \(n_{+}v = 0;\)
\item
  \(\mathfrak{h}v\subset kv\) 且 \(\mathfrak{g}^{\alpha}v = 0\) 对一切 \(\alpha \in \mathrm{B}\) 成立。
\item
  \(\Longrightarrow\) (ii): 假设 \(\mathfrak{b}_{+}v\subset kv\) 。存在 \(\lambda \in \mathfrak{h}^*\) 使 \(v\in V^{\lambda}\) 。设 \(\alpha \in \mathrm{R}_{+}\) 。则 \(\mathfrak{g}^{\alpha}.v\subset V^{\lambda}\cap V^{\lambda +\alpha} = 0\) 。故 \(\mathfrak{n}_{+}v = 0\) 。
\item
  \(\Longrightarrow\) (iii): 这是显然的。
\item
  \(\Longrightarrow\) (i): 这由 \((X_{\alpha})_{\alpha \in \mathrm{B}}\) 生成 \(\mathfrak{n}_{+}\) 这一事实得出（§3 第 3 节 命题 9 (iii)）。
\end{enumerate}

定义 1. 设 V 是一个 g-模，\(v \in V\) 。若 \(v \neq 0\) 且 v 满足引理 1 的条件，则称 v 是 V 的一个本原元。

本原元属于诸 \(V^{\lambda}\) 之一。对一切 \(\lambda\in h^{*}\) ，用 \(V_{\pi}^{\lambda}\) 表示由这样的 \(v\in V^{\lambda}\) 组成的集合：\(b_{+}v\subset kv\) 。于是，权为 \(\lambda\) 的本原元就是 \(V_{\pi}^{\lambda}\) 的非零元素。

命题 1. 设 V 是一个 g-模，v 是 V 的一个本原元，\(\omega\) 是 v 的权。假设 V 作为 g-模由 v 生成。

\begin{enumerate}
\def\labelenumi{(\roman{enumi})}
\item
  若用 \(\mathrm{U}(\mathfrak{n}_{-})\) 表示 \(\mathfrak{n}_{-}\) 的包络代数，则我们有 \(\mathrm{V} = \mathrm{U}(\mathfrak{n}_{-}).v\) 。
\item
  对一切 \(\lambda \in \mathfrak{h}^*\) ，\(V^{\lambda}\) 是由这样的 \(x \in V\) 组成的集合：\(Hx = \lambda(H)x\) 对一切 \(H \in \mathfrak{h}\) 成立。我们有 \(V = \bigoplus_{\lambda \in \mathfrak{h}^*} V^{\lambda}\) ，且每个 \(V^{\lambda}\) 是有限维的。空间 \(V^{\omega}\) 的维数是 1 ，且 \(V\) 的每个权都形如 \(\omega - \sum_{\alpha \in B} n_{\alpha}. \alpha\) ，其中诸 \(n_{\alpha}\) 是 \(\geq 0\) 的整数。
\item
  V 是一个不可分解的 \(\mathfrak{g}\) -模，且它的交换子约化为纯量。
\item
  设 \(\mathrm{U}(\mathfrak{g})\) 是 g 的包络代数，Z 是 \(\mathrm{U}(\mathfrak{g})\) 的中心。存在唯一的从 Z 到 k 的同态 \(\chi\) ，使对一切 \(z \in Z\) ，\(z_{V}\) 是比为 \(\chi(z)\) 的位似。
\end{enumerate}

设 \(\mathrm{U}(\mathfrak{b}_{+})\) 是 \(\mathfrak{b}_{+}\) 的包络代数。我们有 \(\mathrm{U}(\mathfrak{g}) = \mathrm{U}(\mathfrak{n}_{-})\) . \(\mathrm{U}(\mathfrak{b}_{+})\) （第一章 §2 第 7 节 定理 1 的推论 6）。于是

\[
\mathrm{V} = \mathrm{U} (\mathfrak {g}). v = \mathrm{U} (\mathfrak {n} _ {-}). \mathrm{U} (\mathfrak {b} _ {+}). v = \mathrm{U} (\mathfrak {n} _ {-}). v.
\]

用 \(\alpha_{1},\ldots,\alpha_{n}\) 表示 \(R_{+}\) 的不同元素。则

\[
\left(X _ {- \alpha_ {1}} ^ {p _ {1}} X _ {- \alpha_ {2}} ^ {p _ {2}} \dots X _ {- \alpha_ {n}} ^ {p _ {n}}\right) (p _ {1}, \ldots , p _ {n}) \in \mathbf {N} ^ {n}
\]

是 \(\mathrm{U}(\mathfrak{n}_{-})\) 的一个基，故

\[
\mathrm{V} = \sum_ {(p _ {1}, \dots , p _ {n}) \in \mathbf {N} ^ {n}} k X _ {- \alpha_ {1}} ^ {p _ {1}} \dots X _ {- \alpha_ {n}} ^ {p _ {n}} v.\tag{1}
\]

对 \(\lambda \in \mathfrak{h}^*\) ，置

\[
\mathrm{T} _ {\lambda} = \sum_ {(p _ {1}, \ldots , p _ {n}) \in \mathbf {N} ^ {n}, \omega - p _ {1} \alpha_ {1} - \dots - p _ {n} \alpha_ {n} = \lambda} k X _ {- \alpha_ {1}} ^ {p _ {1}} \ldots X _ {- \alpha_ {n}} ^ {p _ {n}} v.
\]

由第七章 §1 第 1 节 命题 2 (ii)，若 \(h \in \mathfrak{h}\) ，则 \(h_{\mathrm{V}}|_{\mathrm{T}_{\lambda}}\) 是比为 \(\lambda(h)\) 的位似。故 \(\mathrm{T}_{\lambda} \subset \mathrm{V}^{\lambda}\) 。另一方面，(1) 蕴含

\[
\mathrm{V} = \sum_ {\lambda \in \omega - \mathbf {N} \alpha_ {1} - \dots - \mathbf {N} \alpha_ {n}} \mathrm{T} _ {\lambda}.
\]

诸 \(V^{\lambda}\) 的和是直和（第七章 §1 第 1 节 命题 3）。由这些观察得出：\(V^{\lambda} = T_{\lambda}\) ；V 是诸 \(V^{\lambda}\) 的直和；且 \(V^{\lambda}\) 是由这样的 \(x \in V\) 组成的集合：\(hx = \lambda(h)x\) 对一切 \(h \in h\) 成立。另一方面，\(\dim V^{\lambda}\) 至多是这样的 \((p_{1}, \ldots, p_{n}) \in \mathbf{N}^{n}\) 组成的集合的基数：它们满足 \(p_{1}\alpha_{1} + \cdots + p_{n}\alpha_{n} = \omega - \lambda\) 。这就证明：\(V^{\lambda} = 0\) （若 \(\omega - \lambda \notin \sum_{\alpha \in B} N\alpha\) ）；\(\dim V^{\omega} = 1\) ；且诸 \(V^{\lambda}\) 都是有限维的。

设 \(c\) 是 V 的交换子的一个元素。对一切 \(h \in \mathfrak{h}\) ，

\[
h c (v) = c h (v) = \omega (h) c (v),
\]

故 \(c(v) \in V^{\omega}\) ；于是存在 \(t \in k\) 使 \(c(v) = tv\) 。现在，对一切 \((p_1, \ldots, p_n) \in \mathbf{N}^n\) ，

\[
c X _ {- \alpha_ {1}} ^ {p _ {1}} \dots X _ {- \alpha_ {n}} ^ {p _ {n}} v = X _ {- \alpha_ {1}} ^ {p _ {1}} \dots X _ {- \alpha_ {n}} ^ {p _ {n}} c v = t X _ {- \alpha_ {1}} ^ {p _ {1}} \dots X _ {- \alpha_ {n}} ^ {p _ {n}} v
\]

故 \(c = t.1\) 。于是，\(V\) 的交换子约化为纯量。这蕴含 (iv) 以及 \(V\) 不可分解这一事实。

定义 2. 命题 1 (iv) 的同态 \(\chi\) 称为 \(\mathfrak{g}\) -模 V 的中心特征标。

命题 2. 设 V 是由权为 \(\omega\) 的本原元 e 生成的 g-模，X 是一个半单 g-模。设 \(\Phi\) 是从 g-模 V 到 g-模 X 的同态组成的集合。则 \(\varphi\mapsto\varphi(e)\) 是从向量空间 \(\Phi\) 到向量空间 \(X_{\pi}^{\omega}\) 的同构。

显然 \(\varphi(e) \in X_{\pi}^{\omega}\) 对一切 \(\varphi \in \Phi\) 成立。若 \(\varphi \in \Phi\) 且 \(\varphi(e) = 0\) ，则 \(\varphi = 0\) ，因为 \(e\) 生成 \(\mathfrak{g}\) -模 V。我们证明：若 \(f\) 是 \(X_{\pi}^{\omega}\) 的非零元素，则存在 \(\varphi \in \Phi\) 使 \(\varphi(e) = f\) 。设 \(X'\) 是 X 的由 \(f\) 生成的子模。由命题 1，\(X'\) 是不可分解的，由于 X 半单，故它是单的。元素 \((e, f)\) 在 \(\mathfrak{g}\) -模 \(V \times X\) 中是本原的。设 N 是 \(V \times X\) 的由 \((e, f)\) 生成的子模。则 \(N \cap X \subset \mathrm{pr}_2(N) = X'\) ，故 \(N \cap X = 0\) 或 \(X'\) ；若 \(N \cap X = X'\) ，则 N 包含线性无关的元素 \((e, f)\) 与 \((0, f)\) ，它们是权为 \(\omega\) 的本原元；这是荒谬的（命题 1），故 \(N \cap X = 0\) 。于是 \(\mathrm{pr}_1|_{N}\) 是从 N 到 V 的单射映射 \(h\) ；这个映射是满射，因为它的像包含 \(e\) 。于是 \(\varphi = \mathrm{pr}_2 \circ h^{-1}\) 是从 \(\mathfrak{g}\) -模 V 到 \(\mathfrak{g}\) -模 X 的、使 \(\varphi(e) = f\) 的同态。

